% Standalone, compile-ready description of the r3 experiment.
% The body can also be copied into a paper or supplement.
\documentclass[10pt]{article}
\usepackage[margin=1in]{geometry}
\usepackage{amsmath}
\usepackage{booktabs}
\usepackage{microtype}
\usepackage[hidelinks]{hyperref}

\newcommand{\CoPE}{\textsc{CoPE}}
\newcommand{\FSRPC}{\textsc{FSR-PC}}
\newcommand{\TBD}{\textit{TBD}}

\title{A Persistent-Lineage, Resource-Bounded Evaluation of \CoPE{} and
\FSRPC{}}
\author{}
\date{}

\begin{document}
\maketitle

\section{Motivation}

In the preceding r2 benchmark, \CoPE{} and \FSRPC{} maintained visibly
different internal task states but achieved identical task-completion rates.
The equality was explained by the execution interface: both representations
were compiled into the same canonical object--target requirements and then
passed through the same pick-and-place stack.  Consequently, lifecycle and
lineage information could improve auditability without affecting behavior.

The r3 experiment removes this insensitivity by constructing a long-horizon
task in which the correct downstream action sequence is bound to a persistent
workflow instance.  The experiment tests the resource-bounded hypothesis that,
under the same model, information, call count, output budget, and wall-clock
budget, persistent typed editing yields higher end-to-end task completion than
repeated complete-state regeneration.

\section{Compared Methods}

At every interruption, both methods receive the same current task state, task
contract, world state, completed progress, current event, full event history,
and compute budget.  Both use one call to the same fixed language model.

\paragraph{\CoPE{}.}
The model emits a minimal typed patch, such as \texttt{Suspend},
\texttt{Override}, \texttt{Restore}, or \texttt{RestoreRoot}.  A persistent
state store applies the operation atomically and owns stable identity,
lifecycle history, parent/root relations, and lineage traversal.

\paragraph{\FSRPC{}.}
The model emits a complete replacement state at every event.  This is a strong
full-state-regeneration baseline: it may preserve stable identifiers and is
explicitly allowed to reconstruct complete lifecycle history and lineage.  It
is not deprived of information.  Its defining constraint is that every slot
and every retained lifecycle record must be serialized again.

After adaptation, both arms use the same plan compiler and the same privileged
LIBERO/MuJoCo geometry-oracle controller.  No learned manipulation policy is
used, so the comparison isolates task-state adaptation rather than visual
perception or low-level control learning.

\section{Lineage-Dependent Task}

Each episode represents a persistent multi-order fulfillment shift.  Most
orders have already been completed but remain referenceable persistent records.
One critical order undergoes suspension, restoration, six nested temporary
overrides, another suspension/restoration cycle, and a final request to withdraw
all detours and resume the depth-zero ancestor.

The original root workflow requires the following final plan:
\begin{enumerate}
    \item place butter in basket B;
    \item place the milk receipt token in basket A;
    \item place the yogurt closure token in basket A.
\end{enumerate}
The deepest detour deliberately has the same first action, butter to basket B,
but binds milk to basket C.  Therefore, the surface object--target pair is
insufficient: selecting the correct remaining actions requires following the
parent/root lineage and recovering the continuation stored by the original
workflow instance.  A registered causal gate confirms that the lineage-aware
plan succeeds while a surface-only plan with the same first action fails exact
task completion.

\section{Fixed Model and Registered Profiles}

The first evaluation fixes the model to \texttt{Qwen/Qwen3-32B}, served locally
through an OpenAI-compatible vLLM endpoint.  Sampling temperature is zero,
top-$p$ is one, and the random seed is fixed to 20260904.  Each event permits
one model call, at most 4096 output tokens, and at most 90 seconds.  Method order
alternates by seed.  Twenty paired seeds are registered for each profile.

\begin{table}[h]
\centering
\caption{Registered r3 profiles.  Token sizes are four-character-per-token
diagnostics used only to check feasibility; server-reported tokenizer usage is
authoritative in the experiment.}
\label{tab:r3-profiles}
\begin{tabular}{lrrrrr}
\toprule
Profile & Initial slots & Lineage depth & Events & Max prompt & Max full state \\
\midrule
Calibration   &  8 & 2 &  7 &  3050 &  2186 \\
Lineage-valid & 10 & 6 & 11 &  5229 &  3765 \\
Long-lineage  & 48 & 6 & 11 & 13398 & 11934 \\
\bottomrule
\end{tabular}
\end{table}

The common model context is 32768 tokens.  Thus, even the estimated
long-lineage prompt plus the registered output allowance remains within the
context window, while an explicit complete-state response is substantially
larger than the 4096-token output allowance.  The small profiles determine
whether the methods agree when both output formats are feasible; the
long-lineage profile is the primary finite-budget condition.

\section{Outcome and Analysis}

The primary endpoint is end-to-end task completion in the long-lineage profile:
\begin{align}
\mathrm{Completion} ={}&
\mathrm{AllEventsProcessed}
\land \mathrm{AllGenerationsValid} \nonumber\\
&\land \mathrm{RootLineageRestored}
\land \mathrm{PlanMatchesRoot}
\land \mathrm{PhysicalSuccess}.
\end{align}
Every attempted episode remains in the denominator.  A timeout, length
truncation, malformed JSON, schema failure, lifecycle-history loss, inconsistent
lineage graph, wrong ancestor, wrong continuation, or physical failure is a task
failure rather than an exclusion.

For each profile, we report completion and valid-generation rates with Wilson
95\% intervals.  The primary paired comparison uses the exact McNemar test and
a paired bootstrap interval for the completion-rate difference
$\CoPE-\FSRPC$.  We additionally report completion among pairs in which both
methods generated valid states.  This conditional quantity is diagnostic: it
can reveal a lineage/continuation error after successful serialization, but it
never replaces the primary endpoint or removes generation failures from it.

\section{Expected Result Pattern}

The preregistered qualitative expectations are:
\begin{enumerate}
    \item \textbf{Calibration:} both methods should usually generate valid
    states and achieve similar completion.  Broad failure of both arms indicates
    a model/prompt integration problem and triggers a stop.
    \item \textbf{Lineage-valid:} both output forms remain feasible.  A \CoPE{}
    advantage here would show lower lifecycle or ancestor-selection error beyond
    serialization length; parity would still be compatible with a strong
    regenerator reconstructing the correct state.
    \item \textbf{Long-lineage:} \CoPE{} should achieve higher primary task
    completion because its patch remains short, whereas complete regeneration
    is exposed to truncation, timeout, and whole-state consistency failures
    under the identical finite budget.
    \item \textbf{Conditional-valid diagnostic:} the gap should be smaller than
    the primary gap.  A remaining gap would support an additional lineage
    correctness benefit; no remaining gap would indicate that the main benefit
    is resource-bounded serialization reliability.
\end{enumerate}

\section{Meaning of the Current Engineering Checks}

With an unlimited deterministic oracle generator, \CoPE{} patches and complete
\FSRPC{} states produce the same correct final-state fingerprint and the same
successful plan.  This is a desirable fairness result.  It shows that the
full-state baseline is capable of representing the correct answer, that both
arms share the same task semantics and executor, and that an eventual bounded
LLM gap cannot be attributed to intentionally making \FSRPC{} semantically
incapable.  It does \emph{not} show that \CoPE{} already has higher empirical
task completion.

Eight of eight local regression tests currently pass.  They cover deterministic
task construction, lineage traversal, exact path equivalence under the oracle,
rejection of lost history, failure accounting for output truncation, the
wrong-but-valid-state diagnostic, paired input fingerprints, and result
analysis.  The separate causal gate also passes, and the packaged handoff was
successfully unpacked and retested.  These checks support implementation
correctness and the validity of the experimental mechanism; they are unit and
integration tests, not eight robot trials and not statistical evidence for the
main hypothesis.

\section{Current Evidential Status}

\begin{table}[h]
\centering
\caption{Evidence available before the fixed-model GPU run.}
\begin{tabular}{p{0.38\linewidth}p{0.16\linewidth}p{0.34\linewidth}}
\toprule
Claim or check & Status & Interpretation \\
\midrule
Lineage is causally load-bearing & Passed & Surface-equivalent first action,
but wrong continuation fails the task. \\
Ideal CoPE/FSR-PC semantic parity & Passed & Removes a straw-baseline and
implementation confound. \\
Local regression suite & 8/8 passed & Supports covered software invariants and
failure accounting. \\
Fixed-LLM generation advantage & \TBD & Requires Qwen3-32B server results. \\
Real LIBERO task-completion advantage & \TBD & Requires the registered 120
episode GPU/server run and controller gate. \\
\bottomrule
\end{tabular}
\end{table}

Accordingly, the current evidence supports the statement that r3 is a valid,
lineage-sensitive, executable test of the hypothesis.  It does not yet support
the empirical conclusion that \CoPE{} has higher task completion than
\FSRPC{}.  That conclusion is warranted only if the preregistered fixed-model
run passes its controller and calibration gates and shows the predicted paired
completion advantage in the long-lineage profile.

\end{document}
